% concatenated from GEA-MEL_hecke_sum_twovar_v3.tex

\documentclass[11pt]{amsart}

\usepackage{amsmath,amssymb,amsthm}

\setlength{\textwidth}{6.5in}
\setlength{\oddsidemargin}{0in}
\setlength{\evensidemargin}{0in}
\setlength{\textheight}{9.0in}
\setlength{\topmargin}{-0.35in}

\newtheorem{theorem}{Theorem}
\newtheorem{lemma}{Lemma}
\newtheorem{corollary}[theorem]{Corollary}
\theoremstyle{remark}
\newtheorem{remark}{Remark}

\newcommand{\Z}{\mathbb Z}
\newcommand{\qbinom}[3][]{\genfrac{[}{]}{0pt}{}{#2}{#3}_{#1}}
\newcommand{\sieve}[2]{\mathcal S_{#1,#2}}

\title{Solutions for Hecke Sum Questions of Banerjee and Bringmann}

\author{George E. Andrews}
\address{The Pennsylvania State University, University Park, Pennsylvania 16802}
\email{andrews@math.psu.edu}

\author{Mohamed El Bachraoui}
\address{Dept. Math. Sci,
United Arab Emirates University, PO Box 15551, Al-Ain, UAE}
\email{melbachraoui@uaeu.ac.ae}

\thanks{First author partially supported by Simons Foundation Grant 633284.}
\date{May 2026}

\begin{document}

\begin{abstract}
The present authors introduced a two-color partition series $S(q)$ and conjectured a Hecke-type formula for the even part of $(q^4;q^4)_\infty S(q)$.  Banerjee and Bringmann proved the conjecture by using indefinite theta functions, modular completions, and Sturm's theorem.  They also asked whether a direct proof, for instance one based on Bailey-type ideas, could be found, and they suggested that the odd residue classes may be worth studying.  We prove a two-variable refinement with an additional parameter $a$.
 Our proof relies entirely on $q$-series combined with the Bailey pairs
The original even identity and the odd identity then follow as corollaries by letting $a=1$.
We also record parameter symmetries and cyclotomic companions, including a vanishing result at $a=i$.
\end{abstract}

\subjclass[2020]{05A17, 11P81, 11P83, 33D15}
\keywords{$q$-series, Bailey pairs, Bailey transform, two-color partitions, Hecke-type sums, finite theta sums, finite pentagonal identities}

\maketitle

\section{Introduction}

We use the standard $q$-series notation as for instance in~\cite[Chapter 1]{GasperRahman}
\[
(a;q)_0:=1,\qquad (a;q)_n:=\prod_{j=0}^{n-1}(1-aq^j),\qquad
(a;q)_\infty:=\prod_{j=0}^{\infty}(1-aq^j),
\]
and
\[
(a_1,\ldots,a_r;q)_n:=(a_1;q)_n\cdots(a_r;q)_n.
\]
For a series $F(q)=\sum_{n\ge0} A(n)q^n$, put
\[
F(q)\big|\sieve{N}{r}:=\sum_{n\equiv r\,({\rm mod}\,N)} A(n)q^n.
\]

The present authors studied the series
\[
S(q):=\sum_{m\ge0}\frac{(q;q^2)_{m}^2q^{2m}}{(-q^2;q^2)_{m}}
\]
in connection with two-color partitions with odd smallest part \cite{AndrewsBachraoui}.  The present authors conjectured a Hecke-type formula for the even part of $(q^4;q^4)_\infty S(q)$.  Banerjee and Bringmann proved this conjecture by completing both sides to modular objects and then using Sturm's theorem \cite{BanerjeeBringmann}.  At the end of their paper they asked whether a more direct proof could be found, for example by using Bailey pairs, and they also suggested studying the odd residue classes.

The main result is a two-variable refinement.  Let $a$ be an indeterminate, or a nonzero complex number, and define
\begin{equation}\label{Sa-def}
S_a(q):=\sum_{m\ge0}\frac{(aq;q^2)_m(q/a;q^2)_m q^{2m}}{(-q^2;q^2)_m},
\qquad
P_a(q):=(q^4;q^4)_\infty S_a(q).
\end{equation}
Thus $S_1(q)=S(q)$.

The following theorem gives the two parity parts of $P_a(q)$.

\begin{theorem}\label{thm:twovar}
For $|q|<1$,
\begin{equation}\label{twovar-even}
P_a(q)\big|\sieve{2}{0}
=
\sum_{n\ge0}(-1)^nq^{6n^2+4n}(1+q^{4n+2})
\sum_{j=-n}^{n}(-1)^j a^{2j}q^{-2j^2},
\end{equation}
and
\begin{equation}\label{twovar-odd}
P_a(q)\big|\sieve{2}{1}
=
\sum_{n\ge0}(-1)^{n+1}q^{6n^2+10n+3}(1+q^{4n+4})
\sum_{j=0}^{n}(-1)^j
\big(a^{2j+1}+a^{-2j-1}\big)q^{-2j(j+1)}.
\end{equation}
\end{theorem}

We first record the basic symmetries in the parameter.

\begin{remark}\label{rem:symmetries}
From \eqref{Sa-def},
\[
S_{a^{-1}}(q)=S_a(q),\qquad P_{a^{-1}}(q)=P_a(q),
\]
and
\[
S_{-a}(q)=S_a(-q),\qquad P_{-a}(q)=P_a(-q).
\]
Thus replacing $a$ by $-a$ leaves the even part unchanged and changes the sign of the odd part.  The specializations $a$ and $a^{-1}$ are identical.
\end{remark}

Putting $a=1$ gives the two identities originally wanted for $S(q)$.  The first is the even Hecke-type identity.

\begin{corollary}\label{cor:even}
For $|q|<1$,
\begin{equation}\label{even-main}
(q^4;q^4)_\infty S(q)\big|\sieve{2}{0}
=
\sum_{n\ge0}(-1)^nq^{6n^2+4n}(1+q^{4n+2})
\sum_{|j|\le n}(-1)^jq^{-2j^2}.
\end{equation}
\end{corollary}

The second is the corresponding odd identity.

\begin{corollary}\label{cor:odd}
For $|q|<1$,
\begin{equation}\label{odd-main}
(q^4;q^4)_\infty S(q)\big|\sieve{2}{1}
=
2\sum_{n\ge0}(-1)^{n+1}q^{6n^2+10n+3}(1+q^{4n+4})
\sum_{j=0}^{n}(-1)^jq^{-2j(j+1)}.
\end{equation}
\end{corollary}

The odd identity separates the two odd residue classes modulo $4$.

\begin{corollary}\label{cor:residue}
We have
\[
(q^4;q^4)_\infty S(q)\big|\sieve{4}{1}=0,
\qquad
(q^4;q^4)_\infty S(q)\big|\sieve{4}{3}
=(q^4;q^4)_\infty S(q)\big|\sieve{2}{1}.
\]
Equivalently, if
\[
S_{r,N}:=S(q)\big|\sieve{N}{r},
\]
then
\begin{equation}\label{S14-S34}
S_{1,4}=0,
\qquad
S_{3,4}=S_{1,2}.
\end{equation}
\end{corollary}

The specialization $a=i$ gives a companion whose odd part vanishes.

\begin{corollary}\label{cor:i}
Let $i^2=-1$, and set
\[
T(q):=\sum_{m\ge0}\frac{(-q^2;q^4)_m q^{2m}}{(-q^2;q^2)_m}.
\]
Then
\begin{equation}\label{i-vanishing}
(q^4;q^4)_\infty T(q)\big|\sieve{2}{1}=0,
\end{equation}
and
\begin{equation}\label{i-companion}
(q^4;q^4)_\infty T(q)
=
\sum_{n\ge0}(-1)^nq^{6n^2+4n}(1+q^{4n+2})
\sum_{j=-n}^{n}q^{-2j^2}.
\end{equation}
\end{corollary}

The primitive sixth-root specialization gives another finite-theta companion.

\begin{corollary}\label{cor:rho}
Let $\rho=e^{\pi i/3}$ and put
\[
C_r:=\rho^r+\rho^{-r}=2\cos(\pi r/3)\qquad(r\in\Z).
\]
Also set
\[
U(q):=\sum_{m\ge0}
\frac{(-q^3;q^6)_m q^{2m}}
{(-q;q^2)_m(-q^2;q^2)_m}.
\]
Then
\begin{equation}\label{rho-even}
(q^4;q^4)_\infty U(q)\big|\sieve{2}{0}
=
\sum_{n\ge0}(-1)^nq^{6n^2+4n}(1+q^{4n+2})
\left(1+\sum_{j=1}^{n}C_jq^{-2j^2}\right),
\end{equation}
and
\begin{equation}\label{rho-odd}
(q^4;q^4)_\infty U(q)\big|\sieve{2}{1}
=
\sum_{n\ge0}(-1)^{n+1}q^{6n^2+10n+3}(1+q^{4n+4})
\sum_{j=0}^{n}C_{j-1}q^{-2j(j+1)}.
\end{equation}
The coefficients $C_j$ are periodic of period $6$:
\[
C_0=2,\quad C_1=1,\quad C_2=-1,\quad C_3=-2,\quad
C_4=-1,\quad C_5=1.
\]
\end{corollary}

The primitive cubic-root specialization gives the complementary cyclotomic companion.

\begin{corollary}\label{cor:omega}
Let $\omega=e^{2\pi i/3}$ and put
\[
D_r:=2\cos(\pi r/3)\qquad(r\in\Z).
\]
Also set
\[
V(q):=\sum_{m\ge0}
\frac{(q^3;q^6)_m q^{2m}}
{(q;q^2)_m(-q^2;q^2)_m}.
\]
Then
\begin{equation}\label{omega-even}
(q^4;q^4)_\infty V(q)\big|\sieve{2}{0}
=
\sum_{n\ge0}(-1)^nq^{6n^2+4n}(1+q^{4n+2})
\left(1+\sum_{j=1}^{n}D_jq^{-2j^2}\right),
\end{equation}
and
\begin{equation}\label{omega-odd}
(q^4;q^4)_\infty V(q)\big|\sieve{2}{1}
=
\sum_{n\ge0}(-1)^{n+1}q^{6n^2+10n+3}(1+q^{4n+4})
\sum_{j=0}^{n}D_{j+2}q^{-2j(j+1)}.
\end{equation}
The coefficients $D_j$ are periodic of period $6$:
\[
D_0=2,\quad D_1=1,\quad D_2=-1,\quad D_3=-2,\quad
D_4=-1,\quad D_5=1.
\]
\end{corollary}

For reference, we recall the form of Bailey's transform that will be used; see Andrews \cite[Chapter 3]{AndrewsQSeries} and Warnaar \cite{WarnaarBailey} for background on Bailey pairs, Bailey's lemma, and conjugate Bailey pairs.  A pair $(\alpha_R,\beta_N)_{R,N\ge0}$ is a Bailey pair relative to $1$, with base $q$, if
\begin{equation}\label{bailey-pair-def}
\beta_N=
\sum_{R=0}^{N}\frac{\alpha_R}{(q;q)_{N-R}(q;q)_{N+R}}.
\end{equation}
A conjugate Bailey pair $(\gamma_R,\delta_N)_{R,N\ge0}$, with the same base, satisfies
\begin{equation}\label{conjugate-bailey-def}
\gamma_R=
\sum_{N=R}^{\infty}\frac{\delta_N}{(q;q)_{N-R}(q;q)_{N+R}}.
\end{equation}
Then Bailey's transform says that
\begin{equation}\label{bailey-transform}
\sum_{N\ge0}\beta_N\delta_N=\sum_{R\ge0}\alpha_R\gamma_R,
\end{equation}
whenever the sums converge.

\section{Two elementary lemmas}

The $q$-binomial coefficient is defined by~\cite[Definition 1]{AndrewsTheory}
\[
\qbinom[q]{M}{K}:=
\begin{cases}
\displaystyle \frac{(q;q)_M}{(q;q)_K(q;q)_{M-K}},&0\le K\le M,\\[6pt]
0,&\text{otherwise}.
\end{cases}
\]

We shall need the finite $q$-binomial theorem~\cite[Eq. (3.3.6)]{AndrewsTheory}
\begin{equation}\label{finite-qbinom}
(z;q)_m=\sum_{j=0}^m(-1)^jq^{j(j-1)/2}\qbinom[q]{m}{j}z^j
\end{equation}
and Euler's expansion~\cite[Eq. (2.2.6)]{AndrewsTheory}
\begin{equation}\label{euler-expansion}
(z;q)_\infty=\sum_{v\ge0}\frac{(-1)^vq^{v(v-1)/2}z^v}{(q;q)_v}.
\end{equation}
In the proof below, Euler's expansion will be used in the special form
\begin{equation}\label{euler-used}
(-zq^{k+1};q)_\infty
=\sum_{v\ge0}\frac{z^vq^{v(k+1)+v(v-1)/2}}{(q;q)_v}.
\end{equation}

The following lemma can be seen as a finite form of Euler's pentagonal theorem.

\begin{lemma}\label{lem:pentagonal}
For $N\ge0$, let
\begin{equation}\label{DN}
D_N(q):=
\sum_{j=0}^{\lfloor N/2\rfloor}
(-1)^jq^{j(3j+1-2N)/2}\qbinom[q]{N-j}{j}.
\end{equation}
Then
\begin{equation}\label{finite-pentagonal}
D_N(q)=
\begin{cases}
(-1)^s q^{-s(3s-1)/2},&N=3s,\\[3pt]
(-1)^s q^{-s(3s+1)/2},&N=3s+1,\\[3pt]
0,&N=3s+2.
\end{cases}
\end{equation}
\end{lemma}

\begin{proof}
The initial values are
\[
D_0(q)=1,\qquad D_1(q)=1,\qquad D_2(q)=0.
\]
We shall prove that, for $N\ge3$,
\begin{equation}\label{DN-recurrence}
D_N(q)=-q^{2-N}D_{N-3}(q).
\end{equation}
The formula \eqref{finite-pentagonal} follows at once by iterating this recurrence.

We use the elementary $q$-Pascal identities~\cite[Eqs. (3.3.3)-(3.3.4)]{AndrewsTheory}
\[
\qbinom[q]{M}{K}=\qbinom[q]{M-1}{K}+q^{M-K}\qbinom[q]{M-1}{K-1}
\]
and
\[
(1-q^K)\qbinom[q]{M}{K}=(1-q^{M-K+1})\qbinom[q]{M}{K-1}.
\]
Applying the first identity to $\qbinom[q]{N-j}{j}$ gives two terms.  In the
second term, the $j=0$ contribution is zero; after replacing $j$ by $j+1$, that
term becomes
\[
-\sum_j(-1)^jq^{j(3j+3-2N)/2}\qbinom[q]{N-2-j}{j}.
\]
Hence
\[
D_N(q)=
\sum_j(-1)^jq^{j(3j+1-2N)/2}\qbinom[q]{N-1-j}{j}
-
\sum_j(-1)^jq^{j(3j+3-2N)/2}\qbinom[q]{N-2-j}{j}.
\]
Apply the first $q$-Pascal identity once more to the first sum.  The part containing
$\qbinom[q]{N-2-j}{j}$ combines with the second sum and gives a factor $1-q^j$.
Using the second displayed identity and then shifting $j$ down by one, we obtain
\[
D_N(q)=
-q^{2-N}
\sum_{j=0}^{\lfloor (N-3)/2\rfloor}
(-1)^jq^{j(3j+1-2(N-3))/2}\qbinom[q]{N-3-j}{j}.
\]
The sum on the right is $D_{N-3}(q)$.  This proves \eqref{DN-recurrence}, and hence the lemma.
\end{proof}

\begin{lemma}\label{lem:tail}
For $|q|<1$,
\begin{equation}\label{tail}
\begin{aligned}
&(q;q)_\infty
\sum_{k\ge0}
\frac{q^k(z^2q^{k+1};q)_k(-zq^{k+1};q)_\infty}{(q;q)_k} \\
&\hspace{35mm}=
\sum_{s\ge0}(-1)^sz^{3s}q^{3s^2+2s}(1+zq^{2s+1}).
\end{aligned}
\end{equation}
\end{lemma}

\begin{proof}
Let the left side of \eqref{tail} be $L(z)$.  By \eqref{finite-qbinom}, with the variable there replaced by $z^2q^{k+1}$, and by Euler's expansion in the form \eqref{euler-used},
\[
(z^2q^{k+1};q)_k
=
\sum_{j=0}^k(-1)^jz^{2j}q^{j(k+1)+j(j-1)/2}\qbinom[q]{k}{j},
\]
and
\[
(-zq^{k+1};q)_\infty
=
\sum_{v\ge0}\frac{z^vq^{v(k+1)+v(v-1)/2}}{(q;q)_v}.
\]
Substitution, followed by the change of variables $k=j+\ell$, gives
\[
L(z)=
\sum_{j,v\ge0}(-1)^jz^{2j+v}
q^{(3j^2+3j+v^2+2jv+v)/2}
\qbinom[q]{j+v}{j}.
\]
Thus
\begin{equation}\label{coeff-L}
[z^N]L(z)=q^{N(N+1)/2}D_N(q).
\end{equation}
By Lemma \ref{lem:pentagonal}, the only nonzero coefficients of $L(z)$ are
\[
[z^{3s}]L(z)=(-1)^sq^{3s^2+2s},
\qquad
[z^{3s+1}]L(z)=(-1)^sq^{3s^2+4s+1}.
\]
This is exactly \eqref{tail}.
\end{proof}

\section{Proof of the two-variable theorem}

We use the following special case of MacMahon's finite form of Jacobi's triple product \cite[p.~75]{MacMahon}:
\begin{equation}\label{macmahon-finite-jacobi}
(zq;q)_M(z^{-1};q)_N
=
\sum_{j=-N}^{M}(-1)^j q^{j(j+1)/2}z^j
\qbinom[q]{M+N}{N+j}.
\end{equation}
Taking $M=N$ and $z=a/q$ in \eqref{macmahon-finite-jacobi} gives
\begin{equation}\label{finite-jacobi}
(a;q)_N(q/a;q)_N
=
\sum_{j=-N}^{N}(-1)^j a^j q^{j(j-1)/2}
\qbinom[q]{2N}{N+j}.
\end{equation}
Now apply \eqref{finite-jacobi} with $q$ replaced by $q^2$ and with $a$ there replaced by $aq$.  This gives
\begin{equation}\label{twovar-finite-jacobi}
(aq;q^2)_N(q/a;q^2)_N
=
\sum_{h=-N}^{N}(-1)^h a^h q^{h^2}
\qbinom[q^2]{2N}{N+h}.
\end{equation}
After division by $(q^2;q^2)_{2N}$ this says that
\begin{equation}\label{beta-a}
\beta_N(a):=\frac{(aq;q^2)_N(q/a;q^2)_N}{(q^2;q^2)_{2N}}
\end{equation}
is a Bailey pair relative to $1$, with $q$ replaced by $q^2$ in \eqref{bailey-pair-def}, where
\begin{equation}\label{alpha-a}
\alpha_0(a)=1,
\qquad
\alpha_R(a)=(-1)^Rq^{R^2}\big(a^R+a^{-R}\big)\quad(R\ge1).
\end{equation}

The conjugate pair used with it is obtained from
\begin{equation}\label{delta-def}
\delta_N=(q^4;q^4)_\infty q^{2N}
\frac{(q^2;q^2)_{2N}}{(-q^2;q^2)_N}.
\end{equation}
For $R\ge0$ set
\begin{align}
\gamma_R
&=\sum_{N=R}^{\infty}
\frac{\delta_N}{(q^2;q^2)_{N-R}(q^2;q^2)_{N+R}} \notag\\
&=q^{2R}(q^2;q^2)_\infty
\sum_{k\ge0}
\frac{q^{2k}(q^{4R+2k+2};q^2)_k(-q^{2R+2k+2};q^2)_\infty}
{(q^2;q^2)_k}. \label{gamma-tail-sum}
\end{align}
Here we used
\[
(q^4;q^4)_\infty=(q^2;q^2)_\infty(-q^2;q^2)_\infty.
\]
Lemma \ref{lem:tail}, with $q$ replaced by $q^2$ and $z=q^{2R}$, gives
\begin{equation}\label{gamma-eval}
\gamma_R=q^{2R}
\sum_{s\ge0}(-1)^sq^{6s^2+(6R+4)s}(1+q^{2R+4s+2}).
\end{equation}
Thus Lemma \ref{lem:tail} is the evaluation of the conjugate Bailey pair.

Since
\[
P_a(q)=\sum_{N\ge0}\beta_N(a)\delta_N,
\]
Bailey's transform \eqref{bailey-transform}, with $q$ replaced by $q^2$, gives
\begin{align}
P_a(q)
&=\sum_{R\ge0}\alpha_R(a)\gamma_R \notag\\
&=\sum_{h\in\Z}(-1)^h a^h q^{h^2}\gamma_{|h|}. \label{master-sum}
\end{align}
This is the basic identity from which both parity parts follow.

Since \eqref{gamma-eval} contains only even powers of $q$, the term $q^{h^2}\gamma_{|h|}$ has the same parity as $h$.  Therefore the even part of \eqref{master-sum} is obtained by putting $h=2r$:
\begin{equation}\label{twovar-even-first}
P_a(q)\big|\sieve{2}{0}
=
\sum_{r\in\Z}a^{2r}q^{4r^2}\gamma_{2|r|}.
\end{equation}
For $r\ge0$, \eqref{gamma-eval} gives
\[
q^{4r^2}\gamma_{2r}
=
q^{4r^2+4r}
\sum_{s\ge0}(-1)^sq^{12rs+6s^2+4s}(1+q^{4r+4s+2}).
\]
Putting $n=r+s$, this becomes
\begin{equation}\label{Er-tail}
q^{4r^2}\gamma_{2r}
=
(-1)^rq^{-2r^2}
\sum_{n\ge r}(-1)^nq^{6n^2+4n}(1+q^{4n+2})
\qquad(r\ge0).
\end{equation}
The same formula with $r$ replaced by $|r|$ holds for all integers $r$.  Substitution in \eqref{twovar-even-first} and interchange of sums give
\[
P_a(q)\big|\sieve{2}{0}
=
\sum_{n\ge0}(-1)^nq^{6n^2+4n}(1+q^{4n+2})
\sum_{|r|\le n}(-1)^r a^{2r}q^{-2r^2},
\]
which is \eqref{twovar-even}.

For the odd part, put $h=\pm(2r+1)$ in \eqref{master-sum}.  Since $(-1)^h=-1$ for odd $h$,
\begin{equation}\label{twovar-odd-first}
P_a(q)\big|\sieve{2}{1}
=
-\sum_{r\ge0}\big(a^{2r+1}+a^{-2r-1}\big)
q^{(2r+1)^2}\gamma_{2r+1}.
\end{equation}
From \eqref{gamma-eval},
\[
q^{(2r+1)^2}\gamma_{2r+1}
=
q^{4r^2+8r+3}
\sum_{s\ge0}(-1)^sq^{6s^2+(12r+10)s}
(1+q^{4r+4s+4}).
\]
Substitution into \eqref{twovar-odd-first} gives
\[
P_a(q)\big|\sieve{2}{1}
=
\sum_{r,s\ge0}(-1)^{s+1}
\big(a^{2r+1}+a^{-2r-1}\big)
q^{4r^2+8r+6s^2+(12r+10)s+3}
(1+q^{4r+4s+4}).
\]
Now put $n=r+s$.  Then
\[
4r^2+8r+6s^2+(12r+10)s
=
6n^2+10n-2r(r+1),
\]
and $(-1)^{s+1}=(-1)^{n+1}(-1)^r$.  Therefore
\[
P_a(q)\big|\sieve{2}{1}
=
\sum_{n\ge0}(-1)^{n+1}q^{6n^2+10n+3}(1+q^{4n+4})
\sum_{r=0}^{n}(-1)^r
\big(a^{2r+1}+a^{-2r-1}\big)q^{-2r(r+1)},
\]
which is \eqref{twovar-odd}.  The theorem is proved.

\begin{proof}
For Corollaries \ref{cor:even} and \ref{cor:odd}, put $a=1$ in Theorem \ref{thm:twovar}.  The even formula becomes \eqref{even-main}, and the odd formula becomes \eqref{odd-main}.
\end{proof}

\begin{proof}
For Corollary \ref{cor:residue}, every exponent on the right side of \eqref{odd-main} is congruent to $3$ modulo $4$, because
\[
3+6n^2+10n-2j(j+1)\equiv3\pmod4,
\]
and the factor $q^{4n+4}$ does not change the residue class.  Hence the $4n+1$ part of $(q^4;q^4)_\infty S(q)$ is zero, and the $4n+3$ part is its whole odd part.  Since $(q^4;q^4)_\infty$ has only powers $q^{4n}$ and has constant term $1$, the same conclusion holds for $S(q)$ itself.  This proves \eqref{S14-S34}.
\end{proof}

\begin{proof}
For Corollary \ref{cor:i}, put $a=i$ in Theorem \ref{thm:twovar}.  Then
\[
(iq;q^2)_m(-iq;q^2)_m=(-q^2;q^4)_m,
\]
so $S_i(q)=T(q)$.  In the even formula, $(-1)^j i^{2j}=1$.  In the odd formula,
\[
i^{2j+1}+i^{-2j-1}=0.
\]
This gives \eqref{i-vanishing} and \eqref{i-companion}.
\end{proof}

\begin{proof}
For Corollary \ref{cor:rho}, put $a=\rho=e^{\pi i/3}$ in Theorem \ref{thm:twovar}.  Since
$\rho+\rho^{-1}=1$,
\[
(\rho q;q^2)_m(\rho^{-1}q;q^2)_m
=\prod_{r=0}^{m-1}(1-q^{2r+1}+q^{4r+2})
=\frac{(-q^3;q^6)_m}{(-q;q^2)_m}.
\]
Thus $S_\rho(q)=U(q)$.  Also, since $\rho^3=-1$,
\[
(-1)^j\rho^{2j}=\rho^{-j}.
\]
Hence the even inner sum in \eqref{twovar-even} becomes
\[
\sum_{j=-n}^{n}\rho^{-j}q^{-2j^2}
=1+\sum_{j=1}^{n}(\rho^j+\rho^{-j})q^{-2j^2}.
\]
For the odd inner sum,
\[
(-1)^j(\rho^{2j+1}+\rho^{-2j-1})
=\rho^{1-j}+\rho^{j-1}=C_{j-1}.
\]
The two stated formulas follow.
\end{proof}

\begin{proof}
For Corollary \ref{cor:omega}, put $a=\omega=e^{2\pi i/3}$ in Theorem \ref{thm:twovar}.  Since
$\omega+\omega^{-1}=-1$,
\[
(\omega q;q^2)_m(\omega^{-1}q;q^2)_m
=\prod_{r=0}^{m-1}(1+q^{2r+1}+q^{4r+2})
=\frac{(q^3;q^6)_m}{(q;q^2)_m}.
\]
Thus $S_\omega(q)=V(q)$.  Put $\rho=e^{\pi i/3}$, so that $\omega=\rho^2$ and $(-1)^j=\rho^{3j}$.  Then
\[
(-1)^j\omega^{2j}=\rho^j.
\]
Hence the even inner sum in \eqref{twovar-even} becomes
\[
\sum_{j=-n}^{n}\rho^jq^{-2j^2}
=1+\sum_{j=1}^{n}(\rho^j+\rho^{-j})q^{-2j^2}
=1+\sum_{j=1}^{n}D_jq^{-2j^2}.
\]
For the odd inner sum,
\[
(-1)^j(\omega^{2j+1}+\omega^{-2j-1})
=\rho^{j+2}+\rho^{-j-2}=D_{j+2}.
\]
The two stated formulas follow.
\end{proof}

\section{Concluding comments}

The proof gives a direct $q$-series answer to Banerjee and Bringmann's question.  MacMahon's finite Jacobi identity supplies the Bailey pair, and Lemma \ref{lem:tail} evaluates the conjugate Bailey pair.  The Bailey transform then gives the two-variable identity in one step.  The specialization $a=1$ gives the original even Hecke-type identity and the odd identity.  The odd identity shows that one odd residue class vanishes and that the other is the whole odd part.

The parameter $a$ records the Jacobi index.  This explains why the even part contains the Laurent powers $a^{2j}$, whereas the odd part contains the paired powers $a^{2j+1}+a^{-2j-1}$.  The symmetry $P_{-a}(q)=P_a(-q)$ shows that replacing $a$ by $-a$ only changes the sign of the odd part.  The case $a=i$ is especially simple: the paired odd weight vanishes, and so the whole odd part is zero.  The specializations $a=e^{\pi i/3}$ and $a=e^{2\pi i/3}$ give two cyclotomic companions, with periodic coefficients $2\cos(\pi j/3)$ in the finite theta sums.


\section*{Declarations}

%\noindent\textbf{Funding.} The first author was partially supported by Simons Foundation Grant 633284.

\noindent\textbf{Conflict of interest.} The authors declare that they have no conflict of interest.

\noindent\textbf{Data availability.} Data sharing is not applicable to this article as no datasets were generated or analyzed during the current study.


\begin{thebibliography}{99}

\bibitem{AndrewsTheory}
G. E. Andrews, \emph{The Theory of Partitions}, Addison-Wesley, Reading, MA, 1976; reprinted by Cambridge University Press, Cambridge, 1998.

\bibitem{AndrewsQSeries}
G. E. Andrews, \emph{$q$-Series: Their Development and Application in Analysis,
Number Theory, Combinatorics, Physics, and Computer Algebra}, CBMS Regional
Conference Series in Mathematics, vol. 66, American Mathematical Society,
Providence, RI, 1986.

\bibitem{AndrewsBachraoui}
G. E. Andrews and M. El Bachraoui, \emph{Congruences for two-color partitions with odd smallest part}, arXiv:2410.14190.

\bibitem{BanerjeeBringmann}
K. Banerjee and K. Bringmann, \emph{Proof of a conjecture of Andrews and Bachraoui on a Hecke sum}, arXiv:2605.10300v1.

\bibitem{GasperRahman}
G. Gasper and M. Rahman, \emph{Basic Hypergeometric Series}, 2nd ed., Encyclopedia of Mathematics and its Applications, vol. 96, Cambridge University Press, Cambridge, 2004.

\bibitem{MacMahon}
P. A. MacMahon, \emph{Combinatory Analysis}, Vol.~II, Cambridge University Press, Cambridge, 1916; reprinted in \emph{Combinatory Analysis, Volumes I and II}, AMS Chelsea Publishing, vol.~137, American Mathematical Society, Providence, RI, 1984.

\bibitem{WarnaarBailey}
S. O. Warnaar, \emph{50 years of Bailey's lemma}, in \emph{Algebraic
Combinatorics and Applications} (G\"ossweinstein, 1999), Springer, Berlin,
2001, pp. 333--347.

\end{thebibliography}

\end{document}
